\documentclass[11pt]{article}
\usepackage[T1]{fontenc}
\usepackage[utf8]{inputenc}
\usepackage{lmodern}
\usepackage[margin=1in]{geometry}
\usepackage{microtype}
\usepackage{array,booktabs,tabularx}
\usepackage{enumitem}
\usepackage{hyperref}
\hypersetup{colorlinks=true,linkcolor=blue!45!black,urlcolor=blue!45!black}
\setlist[itemize]{leftmargin=1.5em,itemsep=0.35em,topsep=0.4em}
\renewcommand{\arraystretch}{1.18}
\title{Revision Memorandum\\\large The Measurement Problem as a Problem of Reference}
\author{Support memorandum}
\date{\today}
\begin{document}
\maketitle

\section*{Purpose}
This memorandum records the exact theoremic claims, non-claims, literature posture, and supporting contents accompanying the manuscript. It is included as a support document rather than as part of the main theorem ladder.

\section*{Exact theorem-level claims now visible in the manuscript}
\begin{enumerate}
\item The bare Everettian package of unitary evolution, entanglement, decoherence, and branch-local record stability is physically non-selective.
\item Branch-local determinacy does not by itself fix diachronic reference.
\item Admissible reference assignments satisfy a strict two-case exhaustion: either they preserve confirmational role up to role-preserving equivalence or they diverge in confirmational role.
\item Divergent admissible assignments make comparative hypothesis ordering reference-relative rather than theory-relative.
\item Memory, anticipation, and agency inherit the same dependence on prior reference stabilization.
\item Interpretation packages fall into an exhaustive four-class selector taxonomy: dynamic selection, ontic privileging, semantic restriction, or discourse deflation.
\item The later admissibility corpus strengthens these results retrospectively but is not a hidden premise of the main theorem ladder.
\end{enumerate}

\section*{Explicit non-claims retained in the final manuscript}
\begin{itemize}
\item The paper does not derive the Born rule, the Hilbert-space formalism, or Tsirelson bounds.
\item It does not propose new quantum dynamics.
\item It does not rank all interpretations by total plausibility.
\item It does not use later admissibility vocabulary as unexplained authority in the main body.
\item It does not deny that discourse-deflationary views can dissolve the problem by refusing the target explanatory burden.
\end{itemize}

\section*{Supporting contents}
\begin{itemize}
\item manuscript in \texttt{main.tex};
\item branching diagram in the Everettian diagnosis section;
\item interpretation matrix in the selector taxonomy section;
\item objections-and-replies section in the main body;
\item audit and dependency appendix;
\item retrospective bridge appendix to the later admissibility corpus;
\item archive folder containing source material and the reconstructed archival transcription.
\end{itemize}

\section*{Hardening addendum (v1 \textrightarrow{} v2)}
The manuscript incorporates the targeted hardening patch-set recorded in \texttt{archive/rewrite\_plan\_hardening\_v2.pdf}. In concrete terms, the manuscript now includes: a sharpened central thesis stated as an exact conditional; an explicit three-level separation between physical, referential, and confirmational questions; a theorem-preview of the two-case exhaustion in the introduction; promotion of the stabilization result to the theorem-level statement "Reference-stabilization necessity"; an explicit anti-ad-hoc clarification attached to semantic infrastructure; and promotion of the decision-theoretic result to the theorem "Decision-theoretic dependence".

\section*{Literature posture}
The bibliography is organized by pressure point rather than by accumulation: Everett/decoherence, self-location and decision theory, collapse and hidden variables, instrumentalism/QBism, and the later admissibility manuscripts. Each cluster earns a local argumentative role.

\section*{Recommended compile target}
The intended Overleaf compile target is \texttt{main.tex}. This yields the manuscript with appendices and references. This memorandum may also be compiled separately if desired.

\end{document}
